% The \rmoon/\Umoon/\gmoon macros in main.tex expand to R_{\leftmoon} etc.; the
% \leftmoon glyph needs a symbol package not loaded in the (fixed) preamble. Redefine
% here (document body, before first use in section 2) with a plain "M" (Moon) subscript.
\renewcommand{\rmoon}{R_{\mathrm{M}}}
\renewcommand{\Umoon}{U_{\mathrm{M}}}
\renewcommand{\gmoon}{g_{\mathrm{M}}}

\section{Introduction}
\label{sec:intro}

The next decade will place several independent positioning, navigation, and timing (PNT)
services around the Moon---NASA's LunaNet and its Lunar Communications Relay and
Navigation Systems, ESA's Moonlight programme and its Lunar Communications and Navigation
Service (LCNS), and their national counterparts \citep{israel2020lunanet,%
fienga2024moonlight,lnis2023}. Each such service realizes not one reference object but
\emph{two}: a spatial reference \emph{frame} (an origin, an orientation, and a scale for
lunar coordinates) and a \emph{timescale} (a lunar coordinate time, its offset from
Terrestrial Time, and its rate). The two are realized by the \emph{same} measurements---%
one-way and two-way ranges between orbiters and surface beacons carry both the geometry
and the clocks---so their datum freedoms are, in principle, coupled. The design of a
lunar reference \emph{system}, now under active study by the IAU and in the Moonlight
lunar-reference-frame and lunar-time work \citep{fienga2024moonlight}, must therefore ask
a joint question: what does a lunar PNT network \emph{fail} to determine about the
combined frame and timescale, and which of those failures are spatial, which temporal,
and which genuinely \emph{coupled}?

This paper answers that question by classifying the null space---the geodetic
\emph{datum defect}, or gauge freedom \citep{baarda1973,teunissen2000}---of the
\emph{coupled} frame$\oplus$timescale Fisher information, over a nine-parameter datum that
carries the spatial Helmert seven of our companion datum-identifiability paper
\citep{baweja2026datum} together with two temporal parameters, a clock offset $\dtau$ and
a clock rate $\dalpha$. It is the capstone of a companion arc: P1 classifies the
\emph{spatial} datum \citep{baweja2026datum}; P2 the cross-provider interoperability floor
\citep{baweja2026interop}; P3 the common-mode error class an integrity monitor cannot see
\citep{baweja2026commonmode}; and this paper (P4) adds \emph{time} and, crucially, the
\emph{coupling}. We are deliberate, throughout, about the status of every statement---%
cited, proven, validated, or modelled---because the naive versions of two of the headline
claims are false, and a careful reviewer will (rightly) reject them.

\paragraph{The SI-second convention, adopted up front.}
We fix the unit convention at the outset: the SI second (Cs) and the $c$-defined metre.
This choice is not cosmetic; it \emph{quarantines a definitional trap}. A pseudorange is
$c\,\Delta t$, so a drift of the \emph{unit of time} $\delta\alpha$ masquerades as a
fractional scale $\delta s$ at coefficient $\approx1$: the two are indistinguishable
\emph{if the unit of time is left free}. This is a choice of units, not an estimation
degeneracy, and adopting the SI second removes it once and for all. Under SI, a
pseudorange is scale-bearing through the fixed $c$, so \textbf{scale is observable}---it
is \emph{not} a datum defect, and we never treat it as a free datum direction. The residual
temporal
freedom $\dalpha$ is then the \emph{estimation-deviation} freedom about the
physically-determined secular rate of lunar coordinate time relative to Terrestrial Time,
not the physical rate itself; we word this carefully wherever it appears.

\paragraph{What we prove.}
The core contribution is a classification theorem for the coupled datum defect
(Section~\ref{sec:theorem}), stated in a \emph{basis-invariant} form. The datum defect is
a property of the null \emph{subspace} of the Fisher information, and we prove that its
spatial/temporal/coupled decomposition is a function of that subspace through the
orthogonal projector $\Pnull=\Umat\Umat^\top$, hence invariant under the
rotation-ambiguous choice of null basis $\Umat\mapsto\Umat\bm Q$. This invariance is the
rigor upgrade of the paper: the tempting per-vector measure ``energy of a null vector in
the temporal block'' is \emph{ill-posed}, because a null basis can be rotated freely. We
then prove (Theorem~\ref{thm:classify}) that a fully internal \emph{free-network} one-way
lunar ranging problem under SI units---mutual common-mode clocks, no external tie---has an
exact datum gauge of \emph{eight}: three translations, three rotations, the clock offset
$\dtau$, and the clock rate $\dalpha$, with scale observable. The joint gauge is therefore
\emph{not} the direct sum of the spatial and temporal gauges; and external ties lift each
temporal defect (Theorem~\ref{thm:lift}), carrying the free-network gauge of eight down to
the defect-$0$ tied realization we validate (Section~\ref{sec:anchor}). A sharp corollary bounds what a Fisher's off-diagonal
block can and cannot certify: a non-block-diagonal Fisher is \emph{necessary but not
sufficient} for a coupled \emph{defect} (Remark~\ref{rem:necessary}).

\paragraph{What we model, and the two couplings.}
The network geometry (beacons and orbits) is \Modelled{}---representative, not flown---and
so are the two coupling \emph{magnitudes} (Section~\ref{sec:couplings}), which are
first-principles engineering bounds cross-checked against the published lunar-clock-rate
band, not certified budgets. \emph{Coupling~1} (range$\leftrightarrow$clock-offset) is
operationally strong but a \emph{near}-defect, not an exact one; its classical root is the
GNSS radial/clock degeneracy, which we cite and do not claim
\citep{sturza1988,brown1992,parkinson1988}. Here we state, and correct, a claim whose
naive form is false: it is \emph{not} true that two-way ranging generically decouples the
lunar frame scale from the timescale offset at the datum level. In a geometrically diverse
multi-beacon multi-epoch one-way network the coupled Fisher is full-rank, the
$\{s,\dtau\}$ marginal is positive-definite, and adding two-way ranges lifts the marginal
by the same factor as adding an equal count of \emph{more} one-way ranges (lift ratio
$\approx1.00$, computed on the modelled network); the strong
$1\,\mathrm{m}\leftrightarrow3.336\,\mathrm{ns}$ degeneracy is a
per-user, poor-local-geometry effect that does not aggregate into a datum-level
near-defect. \emph{Coupling~2} (relativistic rate$\leftrightarrow$frame) is fundamentally
weak: the frame estimation datum perturbs the lunar timescale rate by at most
$\sim\!2\times10^{-17}$, some six orders of magnitude below the
$\pm1.7\times10^{-11}$ reference-\emph{surface} definitional band. The rate is
gauge-\emph{robust} to estimation and gauge-\emph{sensitive} to definition.

\paragraph{What we validate on real data.}
Exactly one claim is externally \Validated{} (Section~\ref{sec:anchor}): that the
coupled-Fisher linear-algebra pipeline---preconditioned assembly, eigendecomposition,
basis-invariant null-space classification, and the $\{s,\dtau\}$ Schur marginal---%
reproduces an \emph{independent} numpy/LAPACK computation to relative and absolute error
below $10^{-3}$, on observation rows whose spatial partials carry the \emph{real} DE440
lunar physical-axis orientation \citep{park2021de440}. The two anchored networks are a
defect-$0$ full-rank network and an exact rate-defect network; both have zero coupled
dimension, so the nonzero-coupling ($\Pst\neq\bm0$) branch of the classifier is
\Proven{} by construction and covered by module unit tests---including a rotated-basis
invariance test---and is \emph{not} claimed as externally validated. We state this scope
precisely, because it is the discipline the paper is built on.

\paragraph{Honest-novelty statement.}
To fix the boundary once: the geodetic datum defect and its null-space/gauge machinery are
classical \citep{baarda1973,teunissen2000}; the range$\leftrightarrow$clock (radial/clock)
degeneracy is classical GNSS \citep{sturza1988,brown1992,parkinson1988}; and the
relativistic dependence of a clock's rate on the gravitational potential it sits in is
standard relativistic geodesy and the IAU relativistic timescale framework
\citep{soffel2003iau,mueller2018clocks}. We claim novelty for none of these. Our
contributions are: (i) the \emph{coupled} frame$\oplus$timescale datum-defect
classification in basis-invariant form, with the exact free-network defect count of eight
under SI and the necessary-but-not-sufficient off-diagonal criterion; (ii) the corrected,
datum-level finding that two-way ranging provides no selective scale$\leftrightarrow$offset
decoupling in a diverse network; (iii) the quantified rate$\leftrightarrow$frame coupling
bound ($\leq\!\sim\!2\times10^{-17}$) and the gauge-robust-to-estimation /
gauge-sensitive-to-definition dichotomy; and (iv) the resulting two design laws for a
lunar reference-system standard.

\paragraph{Honesty scope.}
We separate four kinds of statement and label them at every occurrence. The datum-defect
and gauge-freedom machinery, the radial/clock degeneracy, and the relativistic clock-rate
physics are \emph{cited} prior art. The classification theorem, its basis-invariance, the
exact defect count, the necessary-not-sufficient criterion, and the lift results
(Section~\ref{sec:theorem}) are \Proven{} algebra. The single coupled-Fisher
linear-algebra reproduction on real DE440-orientation rows (Section~\ref{sec:anchor}) is
\Validated{} against an independent numerical oracle. Everything else---the network
geometry and the two coupling \emph{magnitudes} (Section~\ref{sec:couplings})---is
\Modelled{} and tagged \modtag{}. There is no flight data, no certified timescale budget,
no technology-readiness claim above analysis, and no single-number operational rate claim;
every number is regenerated from the open engine by one command
(Section~\ref{sec:availability}).
